\documentclass[11pt]{article}
% preamble.tex --- shared packages, macros, and theorem environments.
% % preamble.tex --- shared packages, macros, and theorem environments.
% % preamble.tex --- shared packages, macros, and theorem environments.
% \input{preamble} after \documentclass in each standalone note.

\usepackage{amsmath, amssymb, amsthm}
\usepackage{mathtools}
\usepackage{geometry}
\usepackage{hyperref}
\usepackage{natbib}
\usepackage{graphicx}
\usepackage{booktabs}
\usepackage{algorithm}
\usepackage{algpseudocode}
\usepackage{xcolor}

\geometry{margin=1in}

\newtheorem{proposition}{Proposition}
\newtheorem{remark}{Remark}

\DeclareMathOperator{\tr}{tr}
\DeclareMathOperator{\diag}{diag}
\newcommand{\R}{\mathbb{R}}
\newcommand{\E}{\mathbb{E}}
\newcommand{\Var}{\mathrm{Var}}
\newcommand{\dd}{\,\mathrm{d}}
\newcommand{\bx}{\mathbf{x}}
\newcommand{\bs}{\mathbf{s}}
\newcommand{\bphi}{\boldsymbol{\phi}}
 after \documentclass in each standalone note.

\usepackage{amsmath, amssymb, amsthm}
\usepackage{mathtools}
\usepackage{geometry}
\usepackage{hyperref}
\usepackage{natbib}
\usepackage{graphicx}
\usepackage{booktabs}
\usepackage{algorithm}
\usepackage{algpseudocode}
\usepackage{xcolor}

\geometry{margin=1in}

\newtheorem{proposition}{Proposition}
\newtheorem{remark}{Remark}

\DeclareMathOperator{\tr}{tr}
\DeclareMathOperator{\diag}{diag}
\newcommand{\R}{\mathbb{R}}
\newcommand{\E}{\mathbb{E}}
\newcommand{\Var}{\mathrm{Var}}
\newcommand{\dd}{\,\mathrm{d}}
\newcommand{\bx}{\mathbf{x}}
\newcommand{\bs}{\mathbf{s}}
\newcommand{\bphi}{\boldsymbol{\phi}}
 after \documentclass in each standalone note.

\usepackage{amsmath, amssymb, amsthm}
\usepackage{mathtools}
\usepackage{geometry}
\usepackage{hyperref}
\usepackage{natbib}
\usepackage{graphicx}
\usepackage{booktabs}
\usepackage{algorithm}
\usepackage{algpseudocode}
\usepackage{xcolor}

\geometry{margin=1in}

\newtheorem{proposition}{Proposition}
\newtheorem{remark}{Remark}

\DeclareMathOperator{\tr}{tr}
\DeclareMathOperator{\diag}{diag}
\newcommand{\R}{\mathbb{R}}
\newcommand{\E}{\mathbb{E}}
\newcommand{\Var}{\mathrm{Var}}
\newcommand{\dd}{\,\mathrm{d}}
\newcommand{\bx}{\mathbf{x}}
\newcommand{\bs}{\mathbf{s}}
\newcommand{\bphi}{\boldsymbol{\phi}}


\title{Foundations: LGCP Intensity Inference\\
via the Mat\'ern SPDE Prior and FEM Discretisation}
\author{Jack Heinzel}
\date{Draft --- \today}

\begin{document}
\maketitle

\begin{abstract}
This note collects the \emph{basis-agnostic} foundations of the pipeline for
inferring the intensity function of an inhomogeneous Poisson process over a
$d$-dimensional domain.  The log-intensity is modelled as a Gaussian random
field whose prior is derived from the stochastic partial differential equation
(SPDE) representation of the Mat\'ern covariance family \citep{lindgren2011}.
We set up the Galerkin weak form and the precision matrix $Q$ in terms of a
generic mass matrix $C$ and stiffness matrix $G$, leaving the concrete choice of
basis to the companion \emph{Universal Function Approximators} note (triangulation,
adaptive hat tree, adaptive constant tree).  We then discretise the likelihood,
specify the joint hyperparameter model $(\mu,\kappa,\sigma)$, and describe NUTS
sampling.  Transdimensional refinement and dimension-robust samplers are covered
in the companion \emph{Sampling Methods} note.  Implementation subtleties
recorded here include the domain-area correction for buffer nodes, an $O(n)$
eigenvalue trick for the log-determinant, and a centred parameterisation that
eliminates Neal's funnel.
\end{abstract}

\tableofcontents
\bigskip

% ============================================================
\section{Problem Setting}
% ============================================================

Let $\mathcal{D} \subset \R^d$ be a bounded domain and let
$\{s_1, \ldots, s_n\} \subset \mathcal{D}$ be a realisation of an
\emph{inhomogeneous Poisson process} with intensity function $\lambda(s) \ge 0$.
The likelihood for the observed point pattern is
\begin{equation}
  \log p(\{s_i\} \mid \lambda) = \sum_{i=1}^n \log \lambda(s_i) - \int_{\mathcal{D}} \lambda(s) \dd s.
  \label{eq:poisson-ll}
\end{equation}
Our goal is Bayesian inference for $\lambda(\cdot)$ given only the $n$
observed locations.  We adopt the \emph{Log-Gaussian Cox Process} (LGCP)
model \citep{moller1998lgcp}, in which
\begin{equation}
  \log \lambda(s) = \mu + x(s), \qquad x \sim \mathcal{GP}(0, \sigma^2 C_\nu(\cdot,\cdot;\kappa)),
  \label{eq:lgcp}
\end{equation}
where $\mu \in \R$ is a global mean (intercept), $x$ is a zero-mean Gaussian
random field with Mat\'ern covariance $C_\nu$ parameterised by range $\kappa$
and marginal standard deviation $\sigma$.

% ============================================================
\section{The Mat\'ern SPDE Prior}
% ============================================================

\subsection{Continuous SPDE}

\citet{lindgren2011} show that a Gaussian random field with Mat\'ern covariance
of smoothness $\nu$ and range $1/\kappa$ is the unique stationary solution of the
SPDE
\begin{equation}
  (\kappa^2 - \Delta)^{\alpha/2} \, x(s) = \mathcal{W}(s), \quad s \in \R^d,
  \label{eq:spde}
\end{equation}
where $\mathcal{W}$ is Gaussian white noise, $\Delta$ is the Laplacian, and
$\alpha = \nu + d/2$.  For $d=2$ the integer case $\alpha=2$ ($\nu=1$) is most
common and is the focus of this work.

This is a major result, and unlocks the computation power of SPDEs while preserving the intuitive structure and powerful 

\subsection{Generalised Anisotropic SPDE}

The stiffness operator can be generalised to allow spatially varying diffusion
\begin{equation}
  (\kappa(s)^2 - \nabla \cdot H(s) \nabla)^{\alpha/2} \, x(s) = \mathcal{W}(s),
  \label{eq:spde-gen}
\end{equation}
where $H(s) \in \R^{d\times d}$ is a positive-definite diffusion tensor.
Setting $H = I$ and $\kappa$ constant recovers \eqref{eq:spde}.

\subsection{Marginal Variance}

For the stationary isotropic case in $d=2$ with $\alpha=2$, the marginal
variance of the field is
\begin{equation}
  \sigma_x^2 = \frac{1}{4\pi\kappa^2}.
  \label{eq:margvar}
\end{equation}
For $\kappa = 6$ this gives $\sigma_x \approx 0.047$, far too small for
intensity fields that vary by one or more log-units.  Including a free scale
parameter $\sigma$ (Section~\ref{sec:hyperparams}) is therefore essential.

% ============================================================
\section{FEM Discretisation}
% ============================================================

This section sets up the finite element machinery in a \emph{basis-agnostic}
way: the Galerkin weak form and the precision matrix $Q$ are written in terms of
a generic mass matrix $C$ and stiffness matrix $G$, independent of how the basis
$\{\phi_j\}$ is constructed.  The concrete closed forms for $C$ and $G$ are
supplied by the companion \emph{Universal Function Approximators} note --- the
adaptive tree hat basis, the adaptive piecewise-constant basis, and the
Delaunay-triangulation baseline each give a different instantiation of the same
matrices.

\subsection{Galerkin Weak Form}
\label{sec:galerkin}

The Galerkin method turns the SPDE into a finite-dimensional linear system
in two steps.

\textbf{Step 1 --- approximate the field.}  Represent $x(s)$ as a linear
combination of hat basis functions:
\begin{equation}
  x(s) \approx \sum_{j=1}^m x_j\,\phi_j(s),
  \label{eq:fem-approx}
\end{equation}
reducing the unknown field to $m$ nodal values $\mathbf{x}\in\R^m$.

\textbf{Step 2 --- enforce the PDE weakly.}  Rather than requiring the SPDE
to hold pointwise (which would be too strong for the piecewise-linear
approximation), we require the residual to be orthogonal to every basis
function $\phi_i$.  For the $\alpha=1$ SPDE
$(\kappa^2-\Delta)x=\mathcal{W}$, this means
\begin{equation}
  \int_\mathcal{D} \phi_i(s)\,\bigl[(\kappa^2-\Delta)x(s)\bigr]\dd s
  = \int_\mathcal{D} \phi_i(s)\,\mathcal{W}(s)\dd s, \quad i=1,\ldots,m.
  \label{eq:galerkin-condition}
\end{equation}
The Laplacian term is handled by \emph{integration by parts}
(Green's first identity):
\begin{equation}
  \int_\mathcal{D} \phi_i\,(-\Delta x)\dd s
  = \int_\mathcal{D} \nabla\phi_i\cdot\nabla x\dd s,
  \label{eq:ibp}
\end{equation}
where the boundary term vanishes.  This step is essential: $\phi_i$ is
piecewise linear and its second derivatives are zero almost everywhere, so the
left side of \eqref{eq:ibp} would be unusable; the right side requires only
first derivatives, which are well-defined on each element of the discretisation
(each dyadic cell of the tree basis, or each simplex of a triangulation).

Substituting \eqref{eq:fem-approx} into the left-hand side of
\eqref{eq:galerkin-condition} and collecting:
\begin{equation}
  \sum_j x_j \Bigl[\kappa^2\!\int\!\phi_i\phi_j\dd s
                  +\int\!\nabla\phi_i\cdot\nabla\phi_j\dd s\Bigr]
  = \sum_j x_j\,[\kappa^2 C_{ij} + G_{ij}]
  = [K\mathbf{x}]_i.
  \label{eq:lhs}
\end{equation}
For the right-hand side, projecting white noise onto $\phi_i$ gives
$z_i = \langle\phi_i,\mathcal{W}\rangle$ with
$\mathrm{Cov}(z_i,z_j)=\langle\phi_i,\phi_j\rangle=C_{ij}$,
so $\mathbf{z}\sim\mathcal{N}(0,C)$.

The Galerkin system is therefore the discrete stochastic equation
\begin{equation}
  K\mathbf{x} = \mathbf{z}, \qquad \mathbf{z}\sim\mathcal{N}(0,C),
  \label{eq:galerkin-system}
\end{equation}
with $K=\kappa^2 C+G$ assembling the two bilinear forms identified in
\eqref{eq:lhs}.  Section~\ref{sec:spde-op} derives the precision matrix from
\eqref{eq:galerkin-system}; the concrete mass and stiffness matrices $C$ and $G$
for each basis are given in the companion \emph{Universal Function Approximators}
note.

\subsection{SPDE Operator and Precision Matrix}
\label{sec:spde-op}

The matrix $K = \kappa^2 C + G$ is the FEM discretisation of the operator
$(\kappa^2 - \Delta)$, as assembled in \eqref{eq:lhs}.

\paragraph{Precision for $\alpha=1$.}
From the Galerkin system \eqref{eq:galerkin-system}, $\mathbf{x}=K^{-1}\mathbf{z}$
with $\mathbf{z}\sim\mathcal{N}(0,C)$, so the covariance of $\mathbf{x}$ is
$K^{-1}CK^{-1}$ ($K$ is symmetric) and the precision is
\begin{equation}
  Q_{\alpha=1} = K C^{-1} K.
  \label{eq:Q1}
\end{equation}

\paragraph{Precision for $\alpha=2$.}
The $\alpha=2$ SPDE $(\kappa^2-\Delta)^2 x = \mathcal{W}$ factors as two
applications of $(\kappa^2-\Delta)$.  Applying the Galerkin discretisation
twice gives the same form $Q = KC^{-1}K$, but $C^{-1}$ is dense (the inverse
of a sparse matrix), making it expensive to work with.  Following
\citet{lindgren2011}, we replace $C$ by the \emph{lumped} mass matrix
$\tilde{C}=\diag(C\mathbf{1})$:
\begin{equation}
  Q = K \, \tilde{C}^{-1} \, K.
  \label{eq:Q}
\end{equation}
Because $\tilde{C}$ is diagonal, $\tilde{C}^{-1}$ costs nothing to apply.
The substitution is a consistent approximation: $\tilde{C}\to C$ as the mesh
refines (the lumped and consistent mass matrices converge).  The resulting $Q$
is sparse --- its nonzero pattern is that of $K^2$ restricted to the mesh
graph --- and positive definite for any $\kappa>0$, making the prior proper.

For general $\alpha\ge 1$,
\begin{equation}
  Q = K\bigl(\tilde{C}^{-1}K\bigr)^{\alpha-1}.
  \label{eq:Q-general}
\end{equation}
% ============================================================
\section{Likelihood Discretisation}
% ============================================================

\subsection{Observation Matrix}

In general the log-intensity at an observation is $\log\lambda(s_i) = \mu +
\bphi(s_i)^T\bx$, where $\bphi(s_i)_j = \phi_j(s_i)$ evaluates every active basis
function at $s_i$.  For the adaptive tree bases this is a tensor-product
evaluation (companion \emph{Universal Function Approximators} note); for the
triangulation baseline the row reduces to the barycentric coordinates of the
containing simplex, as follows.

For observation $s_i$ lying in simplex $\tau(i)$ with vertices
$v_0^{(i)}, \ldots, v_d^{(i)}$, let $\lambda_i = (\lambda_{v_0^{(i)}}, \ldots,
\lambda_{v_d^{(i)}})$ be the nodal values.  By piecewise linear interpolation,
\begin{equation}
  \log\hat\lambda(s_i) = \mu + \sum_{j=0}^{d} b_{ij} \, x_{v_j^{(i)}},
\end{equation}
where $b_{ij} \ge 0$ are the barycentric coordinates of $s_i$ in $\tau(i)$.

The sum term in \eqref{eq:poisson-ll} becomes $\sum_i x_{v_j^{(i)}} b_{ij} = A\bx$
where $A \in \R^{n \times m}$ is the sparse barycentric projection matrix.
Because the $i$-th row of $A$ is the barycentric coordinates of $s_i$, and
since barycentric coordinates sum to one, the log-likelihood can be written as
\begin{equation}
  \sum_{i=1}^n \log\lambda(s_i) = n\mu + \mathbf{d}^T \bx,
  \label{eq:data-term}
\end{equation}
where the \emph{data-accumulation vector} $\mathbf{d} \in \R^m$ is
$d_j = \sum_{i : v_j \in \tau(i)} b_{ij}$, assembled by scattering barycentric
weights to their respective nodes.

\subsection{Integral Approximation and Lumped Mass Quadrature}
\label{sec:integral-approx}

The integral $\int_\mathcal{D} \lambda(s) \dd s$ is approximated by a
nodal quadrature rule.  Expanding $\lambda(s) = \exp(\mu + x(s))$ and
integrating the piecewise linear interpolant gives
\begin{equation}
  \int_\mathcal{D} \exp\!\bigl(\mu + x(s)\bigr) \dd s
  \approx \sum_{j=1}^m c_j \exp(\mu + x_j),
  \qquad c_j = \int_\mathcal{D} \phi_j(s) \dd s = \tilde{C}_{jj}.
  \label{eq:integral-approx}
\end{equation}
The lumped weights $c_j = \tilde{C}_{jj} = (C\mathbf{1})_j$ equal the area
contributed to node $j$ across all incident simplices, specifically
$c_j = \sum_{\tau \ni j} |\tau|/(d+1)$.

\subsection{Domain Correction for Buffer Nodes}
\label{sec:obs-bounds}

When the mesh bounding box $\hat{\mathcal{D}}$ is larger than the observation
domain $\mathcal{D}_{\rm obs}$, the buffer nodes $j \notin \mathcal{D}_{\rm obs}$
have $c_j > 0$ but lie outside the region over which the Poisson likelihood is
defined.  Including them inflates the integral, biasing $\mu$ downward.

The fix is to zero the lumped weights for buffer nodes before forming the integral:
\begin{equation}
  \tilde{c}_j = c_j \cdot \mathbf{1}[s_j \in \mathcal{D}_{\rm obs}].
  \label{eq:obs-bounds}
\end{equation}
With this correction, the likelihood integral is $\sum_j \tilde{c}_j \exp(\mu+x_j)$,
which integrates only over the observation domain.

\subsection{Posterior Normalization Property}
\label{sec:gamma-posterior}

\begin{proposition}
Under a flat prior $p(\mu) \propto 1$ and the LGCP likelihood \eqref{eq:poisson-ll}
with the quadrature \eqref{eq:integral-approx}, the conditional posterior of the
expected total count $N = \sum_j \tilde{c}_j \exp(\mu + x_j)$, given $\bx$, is
$\mathrm{Gamma}(n, 1)$ with mean $n$.
\end{proposition}

\begin{proof}
Let $Z(\bx) = \sum_j \tilde{c}_j e^{x_j}$.  The log-posterior in $\mu$,
$p(\mu \mid \bx, \{s_i\}) \propto \exp(n\mu - Z(\bx) e^\mu)$,
is the kernel of a $\mathrm{Gamma}(n, 1/Z(\bx))$ distribution on
$N = Z(\bx)e^\mu$.  Hence $\E[N \mid \bx] = n$ for any realization of
$\bx$.
\end{proof}

This property provides a useful diagnostic: after sampling, the posterior mean
of $\sum_j \tilde{c}_j \exp(\hat\mu + \hat x_j)$ should equal $n$.  Without
the domain correction~\eqref{eq:obs-bounds}, the guarantee holds for the full
mesh domain $\hat{\mathcal{D}}$ rather than $\mathcal{D}_{\rm obs}$.
% ============================================================
\section{Joint Inference for Hyperparameters}
\label{sec:hyperparams}
% ============================================================

\subsection{Parameterisation and Prior}

We infer $(\mu, \kappa, \sigma, \bx)$ jointly.  The model is
\begin{align}
  \log\kappa &\sim \mathcal{N}(\mu_\kappa, \sigma_\kappa^2), \\
  \log\sigma  &\sim \mathcal{N}(\mu_\sigma, \sigma_\sigma^2), \\
  \bx \mid \kappa, \sigma &\sim \mathcal{N}(0,\; \sigma^2 Q(\kappa)^{-1}), \\
  \{s_i\} \mid \mu, \bx &\sim \text{LGCP with } \log\lambda = \mu + \bx.
  \label{eq:model}
\end{align}

\subsection{Centred Parameterisation to Avoid Neal's Funnel}
\label{sec:centred}

A non-centred parameterisation $\bx = \sigma \tilde\bx$ with
$\tilde\bx \sim \mathcal{N}(0, Q^{-1})$ creates a \emph{Neal's funnel}
\citep{neal2003slice}: the joint distribution
$({\sigma, \tilde\bx})$ and $(2\sigma, \tilde\bx/2)$ produce nearly identical
likelihood values, yet the prior on $\tilde\bx$ \emph{improves} as
$\tilde\bx \to 0$, giving the sampler a free gradient towards $\sigma \to \infty$
with $\tilde\bx \to 0$.  Traces of $\log\sigma$ diverge.

The \emph{centred parameterisation} resolves this by working directly with
$\bx$ in \eqref{eq:model}.  The prior log-density is
\begin{equation}
  \log p(\bx \mid \kappa, \sigma) =
    -\frac{1}{2\sigma^2} \bx^T Q \bx
    + \tfrac{1}{2}\log|Q|
    - m \log\sigma
    + \text{const},
  \label{eq:log-prior-x}
\end{equation}
where the $-m\log\sigma$ term (the Jacobian for the change-of-variables
$\bx \to \tilde\bx = \bx/\sigma$) strongly penalises large $\sigma$.
For $m \approx 130$ nodes, increasing $\sigma$ by a factor of 10 costs
$130 \log 10 \approx 299$ log-units in prior density.

\subsection{Log-Determinant via Generalised Eigenvalues}
\label{sec:logdet}

The log-prior \eqref{eq:log-prior-x} requires $\log|Q(\kappa)| = 2\log|K(\kappa)| - \log|\tilde{C}|$.
Recomputing $\log|K|$ at every MCMC step via Cholesky costs $O(m^3)$ per
iteration.  We reduce this to $O(m)$ via the spectral decomposition.

\begin{proposition}[Eigenvalue log-det trick]
Let $G v_i = \mu_i C v_i$ be the generalised eigenproblem (computed once,
$O(m^3)$).  Then for any scalar $\kappa$,
\begin{equation}
  \log|K(\kappa)| = \log|C| + \sum_{i=1}^m \log(\kappa^2 + \mu_i).
  \label{eq:logdet-trick}
\end{equation}
\end{proposition}

\begin{proof}
Factorise $K = \kappa^2 C + G = C(\kappa^2 I + C^{-1}G)$.  Then
$\log|K| = \log|C| + \log|\kappa^2 I + C^{-1}G|$.
The matrix $C^{-1}G$ has eigenvalues $\mu_i$, so
$\log|\kappa^2 I + C^{-1}G| = \sum_i \log(\kappa^2 + \mu_i)$.
\end{proof}

The eigenpairs $\{(\mu_i, v_i)\}$ are computed once before sampling via
\texttt{scipy.linalg.eigh}; within BlackJAX, \eqref{eq:logdet-trick} reduces
the log-determinant evaluation to a single sum over eigenvalues.

\begin{remark}
This trick requires scalar $\kappa$.  For non-stationary fields where
$\kappa(s)$ varies, the pencil $K(\kappa) = \diag(\kappa^2)C + G$ is no longer
a polynomial in a single scalar, and stochastic Lanczos quadrature (SLQ) would
be needed instead.
\end{remark}

\subsection{Full Log-Posterior}

Combining all terms, the log-posterior to be sampled is
\begin{align}
  \log p(\mu, \log\kappa, \log\sigma, \bx \mid \{s_i\})
  &= \underbrace{n\mu + \mathbf{d}^T \bx - \sum_j \tilde{c}_j e^{\mu+x_j}}_{\text{Poisson likelihood}} \notag \\
  &\quad - \underbrace{\frac{1}{2\sigma^2}\bx^T Q \bx + \frac{1}{2}\log|Q| - m\log\sigma}_{\text{SPDE prior on } \bx}
  \label{eq:log-posterior} \\
  &\quad - \underbrace{\frac{(\log\kappa - \mu_\kappa)^2}{2\sigma_\kappa^2}
         + \frac{(\log\sigma - \mu_\sigma)^2}{2\sigma_\sigma^2}}_{\text{hyperparameter priors}}. \notag
\end{align}
% ============================================================
\section{MCMC Sampling with NUTS}
% ============================================================

\subsection{BlackJAX and Window Adaptation}

We use the No-U-Turn Sampler \citep{hoffman2014nuts} as implemented in
BlackJAX \citep{blackjax2024}, a JAX-native library for Bayesian computation.
JAX's \texttt{jit}-compilation and automatic differentiation compute the
gradient of \eqref{eq:log-posterior} efficiently.

The step size and a diagonal mass matrix are tuned during a warm-up phase using
dual-averaging window adaptation.  After warm-up, samples are drawn and thinned
via a nested \texttt{jax.lax.scan} for memory efficiency.

\subsection{Initialisation}

The sampler is initialised at
\begin{equation}
  \mu^{(0)} = \log\!\left(\frac{n}{\sum_j \tilde{c}_j}\right),
  \quad \bx^{(0)} = \mathbf{0},
  \quad \log\kappa^{(0)} = \mu_\kappa,
  \quad \log\sigma^{(0)} = \mu_\sigma.
  \label{eq:init}
\end{equation}
The initialisation $\mu^{(0)}$ ensures the expected total count under the
prior equals $n$, consistent with Proposition~\ref{sec:gamma-posterior}.

\subsection{Posterior Prediction}

Given posterior samples $\{(\mu^{(s)}, \bx^{(s)})\}_{s=1}^S$, the intensity
at any evaluation point $s^* \in \mathcal{D}$ is obtained by barycentric
interpolation:
\begin{equation}
  \hat\lambda^{(s)}(s^*) = \exp\!\left(\mu^{(s)} + \sum_{j=0}^d b_j(s^*) x^{(s)}_{v_j(s^*)}\right),
\end{equation}
where $(v_0, \ldots, v_d)$ are the vertices of the containing simplex and
$(b_0, \ldots, b_d)$ are the barycentric coordinates of $s^*$.
Points outside the mesh return \texttt{NaN}.
% ============================================================
\section{Implementation Summary}
% ============================================================

The pipeline is split across two Python modules and one Jupyter notebook:

\begin{description}
  \item[\texttt{triangulation.py}] Adaptive kd-tree mesh construction
    (\texttt{AdaptiveKDMesh}) and vectorised FEM assembly
    (\texttt{assemble\_fem\_matrices}, \texttt{precision\_matrix}).
    Supports isotropic and anisotropic diffusion tensors $H$.

  \item[\texttt{inference.py}] Log-posterior construction
    (\texttt{make\_log\_posterior}, \texttt{make\_log\_posterior\_with\_hyperparams}),
    NUTS runner (\texttt{run\_nuts}), and posterior prediction
    (\texttt{predict\_intensity}).

  \item[\texttt{spde\_triangulation.ipynb}] Demonstration on synthetic data from
    a Gaussian mixture Poisson process.  Compares fixed-$\kappa$ and
    joint-hyperparameter models; visualises posterior median intensity,
    point-wise uncertainty, and parameter traces.
\end{description}

% ============================================================
\section{Discussion and Extensions}
% ============================================================

\subsection{Extensions Already Scaffolded}

The \texttt{assemble\_fem\_matrices} function accepts a callable $H(s)$
returning a $(T \times d \times d)$ array of per-simplex diffusion tensors,
enabling \emph{non-stationary} anisotropy.  The companion \emph{Universal
Function Approximators} note covers the stationary anisotropic case in detail;
the key remaining obstacle for both cases is the $O(m^3)$ log-determinant of the
anisotropic operator (discussed there).
Likewise, $\kappa(s)$ can be a callable returning per-simplex values; the
eigenvalue log-det trick then no longer applies.

\subsection{Planned: Non-Gaussian Fields (Level 2)}

The current LGCP assumes a Gaussian field, which may be unsuitable for intensity
functions with sharp discontinuities.  Planned extensions include:
\begin{itemize}
  \item \textbf{$\alpha$-stable SPDE}: replace Gaussian white noise $\mathcal{W}$
    with an $\alpha$-stable L\'evy sheet, yielding a Cauchy or heavy-tailed field;
  \item \textbf{Jump-diffusion SPDE}: add a compound Poisson term to the SPDE,
    allowing finite jumps over infinitesimal regions.
\end{itemize}

Both extensions require replacing the Gaussian prior log-density in
\eqref{eq:log-posterior} with a non-Gaussian counterpart and potentially
switching to a different sampler (e.g., elliptical slice sampling or particle
MCMC).

\subsection{Scalability}

The FEM assembly is fully vectorised (no Python loops over simplices) and runs
in seconds for meshes up to $\sim 10^4$ nodes.  For the isotropic model, the
dominant cost is the one-time eigendecomposition (Section~\ref{sec:logdet}),
$O(m^3)$, taking $\sim 1$\,s for $m \approx 500$; each subsequent NUTS step
costs $O(m)$.  For the anisotropic model, the $O(m^3)$ \texttt{slogdet} is paid
\emph{every leapfrog step} (companion \emph{Universal Function Approximators}
note), limiting the practical mesh size to $m \lesssim 500$.  For larger isotropic meshes one would
switch to a sparse Cholesky factorisation (e.g., via \textsc{CHOLMOD}) and
accept $O(m)$ per-step cost without the precomputed eigenvalues.

\bibliographystyle{plainnat}
\bibliography{refs}

\end{document}
